\documentclass[../candidatura.tex]{subfiles}
\graphicspath{{\subfix{../imagenes/}}}
\begin{document}

\section{Apéndice}

% Se puede incorporar imágenes de obras, links a obras sonoras o audiovisuales, entrevistas, material de archivo, etc.

\subsection{Libro Máquinas sonoras}

\textit{Máquinas sonoras: electrónicas y computacionales} es un libro de apuntes en desarrollo, escrito en coautoría con Matías Serrano \parencite{montoyamoraga_maquinas_2026}. Acompaña los cursos DIS8644 Taller de diseño de máquinas electrónicas y DIS8645 Taller de diseño de máquinas computacionales, que propusimos en 2024 en la Escuela de Diseño de la Facultad de Arquitectura, Arte y Diseño de la Universidad Diego Portales, y que dictamos desde 2025 a estudiantes de pregrado de tercer y cuarto año.

El libro reúne las explicaciones técnicas de las máquinas que esta tesis estudia: relojes, secuenciadores, osciladores, amplificadores, moduladores y filtros, junto con fundamentos de electricidad, componentes básicos, circuitos integrados y conceptos de programación como variables, funciones y clases. Los capítulos de esta tesis remiten al libro para esas explicaciones, y se ocupan de la dimensión crítica y teórica de cada máquina.

El libro es también parte de la práctica de esta investigación: es el material con el que estudiantes construyen sus propias máquinas en aula y para sus prácticas.

Su código fuente es abierto, y la versión más reciente se compila automáticamente y se publica en \url{https://disenoudp.github.io/apuntes-maquinas/}.

% TODO: al entregar la tesis, reemplazar el enlace por una versión congelada del libro (etiqueta de git y DOI de Zenodo)

\subsection{Maqui}

Maqui es un secuenciador digital que se puede usar para controlar otros módulos de sonido.

Este trabajo cuenta con la colaboración de la primera persona practicante que ha tenido la empresa piruetas SpA, el diseñador Andrés Martin. Andrés realizó un trabajo de carcasa y de componentes elementales para la generación de un secuenciador.

\subsection{Menatron}

Menatron es una familia de ROMplers que usan como material fuente la obra de la artista chilena Javiera Mena \footcite{montoyamoraga_menatron_2026}. Está desarrollada como módulos para el ecosistema VCV Rack, con un módulo principal, un módulo con el álbum completo \textit{Esquemas juveniles}, y un módulo por cada canción de ese álbum.

Las gráficas de los módulos fueron realizadas por el diseñador Mateo Arce, y la investigación, el concepto y la programación por mí. Menatron fue presentada en 2026 en el Congreso Javiera Mena, en la Universidad de Santiago de Chile.

\begin{figure}[htbp]
    \centering
    \includegraphics[height=0.35\textheight]{./imagenes/menatron/menatron-00-menatron.pdf}
    \hspace{2em}
    \includegraphics[height=0.35\textheight]{./imagenes/menatron/menatron-01-esquemas-juveniles-album.pdf}
    \caption{Módulos de Menatron para VCV Rack: módulo principal (izquierda) y álbum \textit{Esquemas juveniles} completo (derecha)}
    \label{fig:menatron-principal}
\end{figure}

\begin{figure}[htbp]
    \centering
    \includegraphics[height=0.3\textheight]{./imagenes/menatron/menatron-02-al-siguiente-nivel.pdf}
    \hfill
    \includegraphics[height=0.3\textheight]{./imagenes/menatron/menatron-03-esquemas-juveniles.pdf}
    \hfill
    \includegraphics[height=0.3\textheight]{./imagenes/menatron/menatron-04-como-siempre-sone.pdf}
    \hfill
    \includegraphics[height=0.3\textheight]{./imagenes/menatron/menatron-05-sol-de-invierno.pdf}
    \hfill
    \includegraphics[height=0.3\textheight]{./imagenes/menatron/menatron-06-camara-lenta.pdf}
    \hfill
    \includegraphics[height=0.3\textheight]{./imagenes/menatron/menatron-07-casan-no-puedo-bloquear-lo-que-quiero.pdf}
    \hfill
    \includegraphics[height=0.3\textheight]{./imagenes/menatron/menatron-08-cuando-hablamos.pdf}
    \hfill
    \includegraphics[height=0.3\textheight]{./imagenes/menatron/menatron-09-esta-en-tus-manos.pdf}
    \hfill
    \includegraphics[height=0.3\textheight]{./imagenes/menatron/menatron-10-yo-no-te-pido-la-luna.pdf}
    \hfill
    \includegraphics[height=0.3\textheight]{./imagenes/menatron/menatron-11-perlas.pdf}
    \caption{Módulos de Menatron para cada canción del álbum \textit{Esquemas juveniles}}
    \label{fig:menatron-canciones}
\end{figure}


\subsection{Silenkier}

Silenkier es un sintetizador de tipo sampler, que está cargado con distintos samples de audio extraídos programáticamente de los discursos del actual presidente de Chile, que fueron percibidos como silencios por algoritmos de detección de habla.

\end{document}